\documentclass[11pt]{article}
\usepackage[margin=1in]{geometry}
\usepackage{amsmath,amssymb,amsthm,mathtools}
\usepackage{graphicx}
\usepackage[hidelinks]{hyperref}
\usepackage{microtype}

\newtheorem{theorem}{Theorem}
\newtheorem{lemma}{Lemma}
\newtheorem{remark}{Remark}
\newcommand{\T}{\mathbb T}
\newcommand{\Z}{\mathbb Z}
\newcommand{\FL}{\mathcal F L}
\newcommand{\one}{\mathbf 1}

\title{A noncanonical automorphism at the critical polynomial weight}
\author{Automated mathematics run \texttt{fa\_banach\_001}\\
Model: GPT5.6}
\date{9 August 2026}

\begin{document}
\maketitle

\begin{center}
\fbox{\parbox{0.91\linewidth}{\textbf{Status.}
Claimed full negative resolution (counterexample), likely valid pending human
review.  The construction works for every $1<p<\infty$.}}
\end{center}

\section{Source question}

Kuznetsova and Zadeh study algebra automorphisms of
$\ell^p(\Z,\omega_a)$, where
\[
  \omega_a(n)=\max(1,|n|^a).
\]
After proving the canonical classification in their higher-regularity range,
they state on published page 864, in Section 3 immediately after the relevant
corollary: ``For $a=1$ and $p>1$, the question is open''
\cite[p.~864]{KZ}.  In context, the question is whether every automorphism at
the critical weight $\omega_1$ must still be canonical, namely of one of the
forms
\[
 (Tf)_n=\lambda^n f_n,\qquad (Tf)_n=\lambda^n f_{-n}
 \quad(\lambda\in\T).
\]

\begin{figure}[ht]
\centering
\includegraphics[width=.94\linewidth]{figures/open_problem_crop.png}
\caption{The open statement and its immediate context in \cite[p.~864]{KZ}.}
\end{figure}

\section{Definitions}

Write $\T=\mathbb R/\mathbb Z$ and
$e_n(x)=e^{2\pi i n x}$.  For $1<p<\infty$, let
\[
 \FL^p(\T)=\left\{u:\bigl(\widehat u(n)\bigr)_{n\in\Z}
 \in\ell^p(\Z)\right\},
 \qquad
 \|u\|_{\FL^p}=\|\widehat u\|_{\ell^p},
\]
initially in the sense of periodic distributions.  The Fourier transform
identifies $\ell^p(\Z,\omega_1)$ with the completion $A_p^1$ of the
trigonometric polynomials in the norm
\begin{equation}\label{eq:weighted-norm}
 \|f\|_{A_p^1}^p
 =|\widehat f(0)|^p+
   \sum_{n\ne0}|n\widehat f(n)|^p.
\end{equation}
Moreover, convolution of sequences corresponds to pointwise multiplication
in $A_p^1$.

If $q=p/(p-1)$, H\"older's inequality gives
\begin{equation}\label{eq:embedding}
 \|f\|_\infty
 \le \sum_n|\widehat f(n)|
 \le |\widehat f(0)|+
 \left(2\sum_{n\ge1}n^{-q}\right)^{1/q}
 \left(\sum_{n\ne0}|n\widehat f(n)|^p\right)^{1/p}.
\end{equation}
Thus $A_p^1$ embeds continuously into $C(\T)$.

\section{Idea of the proof}

At order two, the source paper obtains $C^1$ regularity and invokes a
Beurling--Helson rigidity theorem.  At the critical first-order norm, corners
are still allowed.  We exploit precisely that endpoint: a rational
piecewise-affine homeomorphism has constant derivative on each branch, while
the branch cutoffs are bounded-variation multipliers on $\FL^p$.

For a rational slope $r/s$, splitting Fourier frequencies modulo $s$ turns a
localized affine composition into finitely many sparse integer-frequency
reindexings.  Each is followed only by multiplication by a bounded-variation
function.  Such multiplication is bounded on $\FL^p$ because, after one
integration by parts, its coefficient convolution is an average of modulated
discrete Hilbert transforms.  Applying this to the derivative of a composed
trigonometric polynomial controls the critical weighted norm.  The same
argument applies to the rational piecewise-affine inverse.

\section{Two Fourier lemmas}

\begin{lemma}[bounded-variation multipliers]\label{lem:bv}
Let $1<p<\infty$ and let $m$ be a complex-valued function of bounded
variation on $\T$.  Then multiplication by $m$ is bounded on $\FL^p(\T)$.
\end{lemma}

\begin{proof}
Let $\mu=Dm$ be the finite complex measure which is the distributional
derivative of $m$.  With the normalization above, for $k\ne0$,
\[
  2\pi i k\widehat m(k)=\widehat\mu(k).
\]
For a finitely supported sequence $a$, convolution by the nonzero Fourier
coefficients of $m$ can therefore be written as
\[
 \sum_{k\ne0}\widehat m(k)a_{n-k}
 =\frac1{2\pi i}\int_\T
   \left(\sum_{k\ne0}\frac{e^{-2\pi i k t}}{k}a_{n-k}\right)d\mu(t).
\]
For fixed $t$, the operator in parentheses is a modulation conjugate of the
discrete Hilbert transform
\[
  Ha(n)=\sum_{k\ne0}\frac{a_{n-k}}{k}.
\]
The discrete Hilbert transform is bounded on $\ell^p(\Z)$ for
$1<p<\infty$; this is the standard discrete singular-integral theorem (see,
for example, the transference treatment in \cite{Grafakos}).  All modulation
conjugates have the same norm.  Minkowski's integral inequality now bounds
the displayed expression by
$C_p\|\mu\|\|a\|_{\ell^p}$.  The $k=0$ term is
$\widehat m(0)a_n$.  Hence convolution by $\widehat m$, equivalently
multiplication by $m$ on $\FL^p$, is bounded.
\end{proof}

\begin{lemma}[a localized rational affine branch]\label{lem:branch}
Let $I\subset\T$ be an interval, let $r,s$ be positive integers, and let
$b\in\mathbb R$.  On trigonometric polynomials define
\[
  S_{I,r/s,b}f(x)=\one_I(x)f\left(\frac rsx+b\right).
\]
Then $S_{I,r/s,b}$ extends boundedly on $\FL^p(\T)$ for every
$1<p<\infty$.
\end{lemma}

\begin{proof}
It is enough to take $r/s$ in lowest terms.  If
$f(x)=\sum_n a_ne_n(x)$, split $n=sk+j$ with $0\le j<s$.  Directly,
\begin{equation}\label{eq:residue}
 \one_I(x)f\left(\frac rsx+b\right)
 =\sum_{j=0}^{s-1}m_j(x)F_j(x),
\end{equation}
where
\[
 m_j(x)=\one_I(x)e^{2\pi i rjx/s},
 \qquad
 F_j(x)=\sum_{k\in\Z}a_{sk+j}e^{2\pi i(sk+j)b}e_{rk}(x).
\]
The Fourier coefficients of $F_j$ are just the $j$th residue class of
$(a_n)$, inserted at the multiples of $r$ and multiplied by unimodular
numbers.  Consequently
\[
 \|F_j\|_{\FL^p}^p=\sum_k|a_{sk+j}|^p.
\]
Each $m_j$, viewed periodically, has bounded variation: it is smooth on $I$
and zero off $I$, with only finitely many endpoint jumps.  Lemma
\ref{lem:bv}, followed by the finite sum in \eqref{eq:residue}, proves the
claim.
\end{proof}

\section{The counterexample}

\begin{theorem}\label{thm:main}
For every $1<p<\infty$, the convolution Banach algebra
$\ell^p(\Z,\omega_1)$ has a noncanonical algebra automorphism.  More
precisely, extend
\[
 \Phi(x)=
 \begin{cases}
   2x,&0\le x\le \frac13,\\[2mm]
   \dfrac{x+1}{2},&\frac13\le x\le1
 \end{cases}
 \qquad\text{by}\qquad \Phi(x+1)=\Phi(x)+1.
\]
Then $\phi(x+\Z)=\Phi(x)+\Z$ is a non-affine circle homeomorphism, and
composition $C_\phi f=f\circ\phi$ is a bounded invertible algebra
automorphism of $A_p^1$.  Conjugating by the Fourier transform gives the
claimed automorphism of $\ell^p(\Z,\omega_1)$.
\end{theorem}

\begin{proof}
The two formulas for $\Phi$ agree at $x=1/3$, and
$\Phi(0)=0$, $\Phi(1)=1$.  Thus the periodic lift is strictly increasing and
induces an orientation-preserving circle homeomorphism.  Its inverse lift is
\[
 \Psi(y)=
 \begin{cases}
   y/2,&0\le y\le \frac23,\\[1mm]
   2y-1,&\frac23\le y\le1,
 \end{cases}
 \qquad \Psi(y+1)=\Psi(y)+1.
\]

Let $f$ be a trigonometric polynomial and set $g=f\circ\phi$.  The function
$g$ is continuous and piecewise smooth.  Since it is continuous at every
breakpoint, its distributional derivative has no point masses.  The usual
chain rule on the two open branches therefore gives
\begin{equation}\label{eq:derivative}
 g'
 =2\one_{[0,1/3]}(f'\circ(2x))
  +\frac12\one_{[1/3,1]}(f'\circ((x+1)/2)).
\end{equation}
Both summands are instances of Lemma \ref{lem:branch}, with slopes $2$ and
$1/2$.  Hence
\begin{equation}\label{eq:deriv-bound}
  \|g'\|_{\FL^p}\le C_p\|f'\|_{\FL^p}.
\end{equation}
The mean coefficient is controlled using \eqref{eq:embedding}:
\[
 |\widehat g(0)|\le\|g\|_\infty=\|f\|_\infty
 \le C_p'\|f\|_{A_p^1}.
\]
Since $\widehat{f'}(n)=2\pi i n\widehat f(n)$, equations
\eqref{eq:weighted-norm} and \eqref{eq:deriv-bound} show that
\[
  \|C_\phi f\|_{A_p^1}\le C_p''\|f\|_{A_p^1}.
\]

Trigonometric polynomials are dense in $A_p^1$, so $C_\phi$ extends to a
bounded operator on $A_p^1$.  The embedding \eqref{eq:embedding} shows that
convergence in $A_p^1$ implies uniform convergence.  Thus this extension is
still literal pointwise composition by $\phi$, rather than merely an abstract
operator defined by completion.

The inverse lift $\Psi$ is also rational piecewise affine.  Repeating the same
argument proves that $C_{\phi^{-1}}$ is bounded on $A_p^1$.  On polynomials,
and then by density,
\[
 C_{\phi^{-1}}C_\phi=C_\phi C_{\phi^{-1}}=I.
\]
Composition preserves products, so $C_\phi$ is a bounded algebra
automorphism.  Under Fourier transform, multiplication in $A_p^1$ is
convolution in $\ell^p(\Z,\omega_1)$, proving the algebra assertion.

Finally, a canonical automorphism has circle symbol $x\mapsto x+c$ or
$x\mapsto -x+c$.  Our symbol is orientation preserving but has slopes $2$
and $1/2$ on sets of positive measure, so it is neither.  The automorphism is
noncanonical.
\end{proof}

\begin{remark}
The construction actually proves a small general statement: every circle
homeomorphism whose lift and inverse are piecewise affine with finitely many
rational slopes induces an automorphism of $A_p^1$ for
$1<p<\infty$.  Only the explicit two-branch example is needed for the source
question.
\end{remark}

\section{Verification and limitations}

\begin{itemize}
\item The argument covers every $1<p<\infty$.  It does not cover $p=1$,
because the discrete Hilbert transform is not bounded on $\ell^1$; this is
consistent with classical Beurling--Helson rigidity.
\item The derivative formula is first used only for trigonometric
polynomials.  Continuity at the breakpoints removes delta masses, and the
bounded operator is then obtained by completion.  This avoids assuming that
an arbitrary $\FL^p$ derivative is a function.
\item A deterministic finite-Fourier smoke test is included in
\texttt{code/numerical\_smoke\_test.py}.  It checks the explicit formulas and
shows stable sample norm ratios, but it is not part of the proof.
\item Bounded novelty search on 9 August 2026: the run's registry, solution,
attempt, and proof-gap indexes were searched by arXiv id, title, and core
terms.  Exact-phrase searches for the open sentence, DOI/title searches,
author-plus-$a=1$ searches, and searches combining piecewise-linear changes
with weighted Fourier--Lebesgue spaces found the source and general
change-of-variable literature, but no later resolution of this exact
question.  This is evidence, not a guarantee, of novelty.
\item Recommended human check: audit Lemmas \ref{lem:bv} and
\ref{lem:branch}, especially the periodic bounded-variation interpretation of
$\one_I e^{2\pi i\alpha x}$, and then the completion argument in Theorem
\ref{thm:main}.
\end{itemize}

\begin{thebibliography}{9}

\bibitem{KZ}
Y. Kuznetsova and S. Zadeh,
\emph{On isomorphisms between weighted $L^p$-algebras},
Canadian Mathematical Bulletin \textbf{64} (2021), no.~4, 853--866;
arXiv:1803.06083; doi:10.4153/S0008439520000880.

\bibitem{LO}
V. Lebedev and A. Olevskii,
\emph{$C^1$ changes of variable: Beurling--Helson type theorem and
H\"ormander conjecture on Fourier multipliers},
Geometric and Functional Analysis \textbf{4} (1994), no.~2, 213--235.

\bibitem{Grafakos}
L. Grafakos,
\emph{Classical Fourier Analysis}, 3rd ed., Graduate Texts in Mathematics
249, Springer, 2014.  See the treatment of singular integrals and
transference underlying the boundedness of the discrete Hilbert transform on
$\ell^p(\Z)$, $1<p<\infty$.

\end{thebibliography}

\end{document}

